\chapter{Experimental Reproducibility and Integration}
\label{chap:reproducibility}
% FIIT: documentation, containerisation, open-science practices; ~3 pages.

All code, firmware and processed data are in a single public repository.\footnote{\url{https://github.com/LubimPiskoty/Correlation-Power-Analysis-of-Neural-Network-Weight-Parameters}}
The pipeline has four stages (firmware, capture, preprocessing, attack), each run by one command.

\section{Hardware and Software Requirements}

\paragraph{Hardware.}
A ChipWhisperer-Lite (CWLITEARM) with its STM32F3 target, connected over USB.
Only the capture stage needs the hardware; the analysis runs on any computer.

\paragraph{Software.}
\begin{itemize}
	\item Python~3 with the packages pinned in \texttt{requirements.txt} (\texttt{chipwhisperer} 6.0.0, \texttt{numpy}, \texttt{scipy}, \texttt{matplotlib}, \texttt{tqdm}, \texttt{jupyter}).
	\item The \texttt{arm-none-eabi-gcc} toolchain and \texttt{make}.
	\item The ChipWhisperer firmware tree, expected at \texttt{\$HOME/chipwhisperer/firmware/mcu} (variable \texttt{FIRMWAREPATH} in the firmware Makefile).
\end{itemize}

\section{Repository Structure}

\begin{description}
	\item[\texttt{firmware/}] C code for the target: the network (\texttt{network.c}), its ground-truth weights (\texttt{network\_config.h}) and the SimpleSerial commands (\texttt{main.c}).
	\item[\texttt{capture/}] \texttt{capture.py} and its settings in \texttt{configs/cw\_config.py}.
	\item[\texttt{analysis/}] \texttt{preprocess.py} and the Jupyter notebooks of the attack.
	\item[\texttt{data/}] \texttt{raw/} (ChipWhisperer projects) and \texttt{processed/} (clipped \texttt{.npz} files).
	\item[\texttt{results/}] Figures and tables generated by the notebooks.
	\item[\texttt{thesis/}, \texttt{paper/}] \LaTeX{} sources.
\end{description}

All capture parameters (number of traces, input range, batch size, seed) are in \texttt{cw\_config.py}.
The weight range is set by \texttt{RANDOM\_WEIGHT\_MAX} in \texttt{main.c} and must match the profiling setting.

\section{Reproducing the Experiments}

\subsection{Building and Flashing the Firmware}

Two firmware builds exist, both with the same code layout.
The profiling build generates random weights from a seed; the attack build keeps the fixed unknown weights.
\begin{minted}{bash}
cd firmware
# profiling
make PLATFORM=CWLITEARM CRYPTO_TARGET=NONE
# attack
make PLATFORM=CWLITEARM CRYPTO_TARGET=NONE IS_ATTACK_MODE=1
\end{minted}
The capture script runs both builds and flashes them itself, so this step is only needed for manual use.

\subsection{Capturing Traces}

\begin{minted}{bash}
python -m venv .venv && source .venv/bin/activate
pip install -r requirements.txt
cd capture && python capture.py
\end{minted}
The script captures 50\,000 profiling and 5\,000 attack traces with seed 148, using a new random weight set every 100 traces.
It writes the ChipWhisperer projects \texttt{data/raw/profile} and \texttt{data/raw/attack}.

\subsection{Preprocessing}

\begin{minted}{bash}
python analysis/preprocess.py
\end{minted}
The script clips the raw projects to the analysed window (samples 0--1000) and saves \texttt{data/processed/profile.npz} and \texttt{attack.npz}.
It must be rerun after every new capture.

\subsection{Running the Attack}

Run \texttt{analysis/TemplateDPA.ipynb} from top to bottom.
It selects points of interest, builds the templates, searches for the weight and computes the success rate against the number of traces.
Figures are written to \texttt{results/figures/}.
There is no test suite; an attack is successful when the recovered weight equals the ground truth in \texttt{firmware/network\_config.h}.

\section{Data Availability and Open Science}

The source code, firmware and the processed traces (\texttt{data/processed/}) are in the public repository, which is enough to rerun the whole analysis without hardware.
The raw ChipWhisperer projects are large and excluded from version control.
Captures are seeded, but they depend on the physical device, so a new capture reproduces the results statistically rather than trace for trace.
Containerisation was not used, because a container cannot easily access the USB-connected ChipWhisperer; the pinned \texttt{requirements.txt} fixes the Python environment instead.
